Let \(g\) be the displayed finite wavelet sum and, for each \(z\), let \(u_z=\mathcal G_{n,z}g(\cdot,z)\). The Neumann boundary conditions and the zero outcome integral of every retained feature give
\[
\langle g,u_z\rangle_{L^2(dy)}
=\int_0^1\bar q_n(y,z)|\partial_yu_z(y)|^2dy.          \tag{1}
\]
Writing \(H_g(y,z)=\int_0^yg(x,z)dx\) and integrating the differential equation gives \(\bar q_n\partial_yu_z=-H_g\). Thus
\[
\langle g,\mathcal G_{n,z}g\rangle
=\int_0^1\frac{H_g(y,z)^2}{\bar q_n(y,z)}dy
\asymp \|g(\cdot,z)\|_{H^{-1}(0,1)}^2.               \tag{2}
\]
The constants depend only on \(c,C\). For the declared boundary wavelets, applying the antiderivative to an outcome wavelet at level \(j_y\) multiplies its \(L^2\) scale by \(2^{-j_y}\); distinct levels have uniformly bounded overlap and the boundary scaling functions form a fixed finite-dimensional block. Hence
\[
\|g\|_{L^2(dR_n;H^{-1}_y)}^2
\asymp\sum_{j_y,k_y,\gamma}2^{-2j_y}x_{j_y,k_y,\gamma}^2. \tag{3}
\]
The bounds \(c\le r_n\le C\) transfer the Lebesgue tensor estimate to \(dR_n\). Combining (1)--(3) proves the first display.

Write \(p_{a,n}=q_{a,n}r_n\) for \(d\ge1\). If \(x\) is the coefficient vector of \(g\), then
\[
x^\top\Gamma_{a,J_y,J_z,n}x
=\inf_{v\in\mathbb R}\int(g-v)^2p_{a,n}\,dy\,dz.       \tag{4}
\]
Because \(\int_0^1g(y,z)dy=0\) for every \(z\),
\[
\int(g-v)^2dy\,dz=\int g^2dy\,dz+v^2.
\]
The two density lower bounds and Parseval therefore give a fixed positive lower bound in (4), while the upper density bounds give a fixed upper bound. For \(d=0\), the same argument uses only \(c\le q_{a,n}\le C\). Summing the two arms proves \(c_\Gamma I\preceq\Gamma\preceq C_\Gamma I\).

Congruence by \(\Gamma^{1/2}\), (3), and the variational characterization of ordered eigenvalues compare the eigenvalues of \(K\) with those of the diagonal reference carrying weight \(2^{-2j_y}\). There are \(\asymp2^{j_y+dJ_z}\) coordinates at outcome level \(j_y\). Thus, for \(2^{-2J_y}\le u\le1\),
\[
N(u)\asymp 2^{dJ_z}\sum_{j_y\le(1/2)\log_2(u^{-1})}2^{j_y}
\asymp\min\{D_{J_y,J_z},2^{dJ_z}u^{-1/2}\}.           \tag{5}
\]
Inverting (5) gives \(\nu_k\asymp\min\{1,(2^{dJ_z}/k)^2\}\). Hence
\[
\sum_k\nu_k^2\asymp
2^{dJ_z}+2^{4dJ_z}\sum_{k>2^{dJ_z}}k^{-4}
\asymp2^{dJ_z},                                       \tag{6}
\]
while the largest eigenvalue is comparable to one. Equations (5)--(6) prove all three spectral assertions and the normalized maximal-eigenvalue bound. When \(d=0\), the same sum has no growing multiplicity and converges to a finite positive constant.

Finally, anisotropic H\"older approximation and the extra antiderivative weight in the outcome coordinate give
\[
\|(I-P_{J_y})\Delta_n\|_{L^2(dR_n;H_y^{-1})}
\lesssim2^{-(s_y+1)J_y},\qquad
\|(I-P_{J_z})\Delta_n\|_{L^2(dR_n;H_y^{-1})}
\lesssim2^{-s_zJ_z}.                                  \tag{7}
\]
The first-order terms have been removed by the correction, so the two deterministic remainders are bounded by the squares of the two norms in (7). Consequently
\[
|B^y_{J_y,n}|\lesssim2^{-2(s_y+1)J_y},
\qquad |B^z_{J_z,n}|\lesssim2^{-2s_zJ_z},
\]
uniformly over the fixed-radius ball.